\section{Two families of simple curves}\label{sec:curves}
We fix the based-word convention, since the order of the factors matters
both topologically and in the tangent calculation. A genus-two surface
with one boundary component is obtained by thickening a rose with cyclic
half-edge order
\begin{equation}\label{eq:rose}
 a,\ b^{-1},\ a^{-1},\ b,\ c,\ d_0^{-1},\ c^{-1},\ d_0.
\end{equation}
Composing the cyclic permutation with edge reversal gives one boundary
cycle with word $[a,b][c,d_0]$. In genus $g$ append the blocks
$a_i,b_i^{-1},a_i^{-1},b_i$, and cap the boundary to obtain $\Sigma_g$.
A turn from a letter $x$ to $y$ is drawn as a chord joining half-edges
$x^{-1}$ and $y$ in the vertex disk.

\begin{lemma}\label{lem:curves}
For all $k,l\geq1$, the words
\begin{align}
 u_0(k,l)&=b a^k d_0 c^l,\label{eq:u0}\\
 u_1(k,l)&=a b a^{k-1}c d_0 c^{l-1}\label{eq:u1}
\end{align}
represent simple nonseparating curves. For every $i\geq3$, so do
$u_1(k,l)a_i$ and $u_1(k,l)b_i^{-1}$.
\end{lemma}
\begin{proof}
First consider the words $bd_0$ and $abcd_0$. Their turn chords in
\eqref{eq:rose}, numbered from one to eight, are respectively
\[
 \{2,8\},\{4,6\};\qquad
 \{3,4\},\{2,5\},\{7,8\},\{1,6\}.
\]
In each case the chords are noninterleaving. Each edge is used once,
so these chords and the edge bands smooth to a single embedded closed
curve. Its homology class is nonzero and primitive, hence it is
nonseparating after capping the boundary.

The oriented Dehn twist about $a$, or its inverse according to the choice
of a positive twist, has based action $b\mapsto ba$ and fixes $a$ and the
other generators. This follows by cutting the $a$-annulus and following the
single transverse $b$-arc across one twist. Iteration gives
$b\mapsto ba^m$. The identical construction on the other handle gives
$d_0\mapsto d_0c^n$. These are actual boundary-fixing twist powers;
algebraically $[a,ba^m]=[a,b]$ and $[c,d_0c^n]=[c,d_0]$, so the chosen
based boundary word is fixed exactly. Apply exponents $(k,l)$ to $bd_0$
and $(k-1,l-1)$ to $abcd_0$. This gives \eqref{eq:u0}--\eqref{eq:u1}
and preserves simplicity and nonseparation.

For an additional handle let $p$ be the position of $a_i$ in the full
cyclic order. The closing chord $\{1,6\}$ of $abcd_0$ is replaced, for
$abcd_0a_i$, by $\{6,p\},\{1,p+2\}$, and for $abcd_0b_i^{-1}$ by
$\{6,p+1\},\{1,p+3\}$. These chords are again noninterleaving with each
other and with the three unchanged chords. Each edge is used once.
The resulting curves are simple and have nonzero homology classes.
The same boundary-fixing twists give the claimed words.
\end{proof}

We will also use the twist about $b$ whose based action is $a\mapsto ab$.
It fixes all other generators and satisfies $[ab,b]=[a,b]$. All twists
used below therefore lie in the orientation-preserving mapping class group
of the closed surface. No replacement of $ba^k$ by $a^kb$ is implicit.
